\subsection{Fault Location, Isolation \& Service Restoration}
\label{sec:flisr}

With resiliency of the network a key concern to the DSO and frequently mentioned as a requirement for the smart grid in the literature, FLISR being one such technique which has been
explored and discussed in various research. \textit{Agüero} \cite{6344960} describes FLISR as a self-healing technique, though does make mention of certain authors preferring the term ``self-restoration'',
as in reality no system can currently fix physical damage to the grid without direct intervention by DSO personnel.
The paper discusses how \ac{da} technologies can enable remote monitoring, coordination and operation of distributed network assets at substation, feeder and customer level.
Most relevant, is the description of utilising remote, on-pole switching devices as a method of sectionalising and restoring supply to customers in the event of a fault.
The paper describes how the need for manual investigation and operating of switching devices on the part of DSO operation staff can be reduced or rendered redundant through automated FLISR,
and the benefits of such a solution being the improvement of reliability statistics and the savings made due to the reduced cost of outages.
\\
However, the text posits that leveraging such devices for self-healing is more appropriate for urban and suburban areas,
rather than remote rural locations, where radial distribution is more commonplace (i.e., one source of power for consumers on the network). Though it could be argued that resilience strategies such as FLISR are more
necessary for rural areas, due to the increased prevalence of faults as a result of deterioration and damage in remote locations.
\\
\\
A report produced by the US Department of Energy in 2014 \cite{us_dept_energy}  aimed to gather metrics on five separate FLISR solutions implemented throughout the United States over the course of one year.
The findings included a large reduction in the number of customers affected by service loss (270,000) with 38 million fewer CML in total, translating to a reduction of 45\% in the amount of interruptions and 51\% in CML per fault event.
The report specifies that an essential component for successful FLISR operations are the communications networks used to monitor and enable control over switching systems and devices, in combination with suitable interfaces which allow
interoperability between various DSO systems and device types on the grid.
\\
Each FLISR system explored utilises various device types, communication methods, modes of operation and location of the deployed system
as part of its implementation. For example, a utility named CenterPoint trialed their FLISR system, named ``Self-Healing Grid'' which is supported by \ac{igsd} and requires manual validation by a
System Operator in order to enact control operations. Conversely, a utility named Duke leveraged electronic reclosers, circuit breakers and line sensors for its FLISR system, dubbed ``Self-Healing Teams'' to enable a fully-automated solution.
\\
The report specifies a number of learning's gathered by the utilities from the FLISR trials, such as the need for the telecommunications network to be more resilient than the power delivery systems which rely upon them, and the necessity for
the utility to consider the additional steps required to successfully deploy and operate a FLISR system. Steps identified included the testing and installation of new equipment, and the frequent firmware and software updates necessary to ensure the system continues to operate
effectively due to its autonomous nature, though one utility, Georgia Power, found that developing a simulator to test and validate FLISR eliminated the need for extensive field trials and served as a training tool for operation staff.
\\
\\
A paper by \textit{Eriksson et al.}\ \cite{7000578} further discusses the limitations and considerations when implementing FLISR solutions, specifically centralised deployments, such as having a single point of failure and the requirement for the handling of large amounts of real-time data, large network topologies and the processing necessary to ensure consistent operation. They propose a
distributed FLISR utilising a \ac{mas} devices, tested against a radial network commonly found in remote areas, and applying Prim's Algorithm to devise the FLISR logic. The solution proposed was tested against a
cyber-physical testbed, with the results finding that the \acs{mas} approach to FLISR provided efficient grid stabalisation and fault restoration.
\\
\\
Most relevant is a report produced by ESB Networks as an output of the H2020 project SOGNO \cite{sogno_report}.\ This project focused on delivering turn-key grid services for power quality monitoring, estimating network state,
load prediction and forecasting as well as FLISR. The FLISR implementation described in the report leveraged communication from Load Break Fault Make (LBFM) on-pole circuit breakers in the field. In the event of a fault,
these switches communicate a Fault Passage Indicator (FPI) and Loss of Voltage Indicator (LVI) and in this way it is possible to identify a fault occurrence. ESB Networks notes that a number of their networks currently rely upon
manual intervention on the part of control room operation staff to manually identify and resolve fault events when they occur. This need for human intervention greatly increases the time taken to restore supply to consumers, leading
to higher CML per event. While the FLISR solution tested as part of the SOGNO project produced good results, ESB Networks noted that its reliance on a centralised instance and telecommunications to operate is a severe limitation.

% subsection flisr (end)
